\section{Discusión}

El recorrido completo deja ver que cada etapa del flujo responde una pregunta distinta, y que
ninguna sustituye a las demás.

El modelo obtenido del balance de energía resultó útil para entender el sistema pero no para
diseñar. Su linealización es válida solo en un entorno reducido del punto de operación, como
cuantifica la Tabla~\ref{tab:fisico}, mientras que el transitorio que debe manejar el controlador
es precisamente un calentamiento desde temperatura ambiente. Por eso el modelo de diseño tuvo que
obtenerse de forma experimental. El FOPDT identificado con el ensayo escalón cumple
esa función, pero la validación cruzada mostró también su límite: los residuos no son ruido
alrededor de cero sino un patrón que crece con el nivel de potencia. Esto significa que el modelo no
falla al azar, sino que subestima la temperatura de forma creciente a medida que se aleja de la
región donde fue identificado, y que un solo modelo lineal con retardo no describe la planta en todo
el rango. La elección del modelo nominal fue, en ese sentido, una elección de región de operación
más que de calidad absoluta.

El efecto de la saturación resultó ser el factor de diseño más relevante de toda la práctica. El PI
sintonizado con SIMC demanda un \SI{274}{\percent} de potencia en el primer instante, casi tres
veces lo que el actuador puede entregar. Con el actuador limitado pero sin protección del
integrador, el controlador sintonizado correctamente incumple el requisito de sobreimpulso por un
factor de seis. La causa no está en la sintonización sino en que el integrador sigue acumulando
mientras la potencia ya no puede crecer, de modo que el estado del controlador deja de corresponder
a lo que realmente ocurre en la planta. El \emph{clamping} corrige eso deteniendo la acumulación
durante la saturación, y el costo es de apenas \SI{10}{\second} de tiempo de subida. La lección es
que el desempeño de un controlador no puede evaluarse sin las restricciones del actuador: el caso
ideal de la Tabla~\ref{tab:mil} cumple los requisitos y a la vez es irrealizable.

La comparación entre los dos diseños refuerza esa idea. Sobre el modelo lineal sin restricciones el
PI--CSD parece claramente superior, con mejor tiempo de establecimiento, menor sobreimpulso y
márgenes mayores. Al incorporar el límite de potencia el orden cambia parcialmente: el PI--SIMC
resulta más rápido porque aprovecha toda la autoridad disponible, mientras que el PI--CSD se
mantiene mucho más cerca de la capacidad real del actuador. Ambos cumplen, y decidir entre ellos
exige un criterio explícito ---rapidez frente a robustez y esfuerzo de control--- que no se deduce
de las métricas por sí solas.

Las verificaciones SIL y PIL cumplieron un papel distinto del de MIL. MIL evalúa si el algoritmo
resuelve el problema de control; SIL y PIL no evalúan desempeño en absoluto, sino conservación de
comportamiento al cambiar de representación. Que ambas dieran igualdad exacta y no solo equivalencia
dentro de una tolerancia da confianza en que las diferencias observadas después, sobre la planta
real, no provienen del código ni del procesador. Ese es su valor práctico: al llegar a la
implementación física, el conjunto de causas posibles de una discrepancia ya está acotado al modelo
de la planta y al hardware de entrada y salida. Al mismo tiempo, conviene no sobreinterpretar el
resultado: durante PIL la planta, el retardo y la temperatura ambiente siguen siendo señales de
Simulink, y no se ejercita ningún canal físico.

Efectivamente, las diferencias aparecieron al cerrar el lazo sobre la TCLab, y son del tipo que las
etapas anteriores no podían revelar. La más significativa es el error de ganancia: la planta
necesita en régimen más potencia que la prevista (Tabla~\ref{tab:experimental}), lo que equivale a
una ganancia estática aparente cerca de un \SI{15}{\percent} menor que la nominal. Esa
diferencia es coherente con el sesgo detectado en la validación cruzada y con la dependencia de la
ganancia respecto de la región de operación, y se traduce en un lazo más lento que el simulado. A
ello se suman efectos ausentes del modelo: el ruido de medición amplificado por
$\Kp = \SI{7.62}{\percent\per\celsius}$, la cuantización del convertidor, y variaciones ambientales
durante los \SI{607}{\second} del experimento. Ninguno de estos efectos invalida el diseño ---el
controlador cumple los requisitos sobre la planta real--- pero sí muestran que el margen con que se
cumplen es menor que el que sugería la simulación, y que ese margen proviene en buena medida de
haber elegido un $\tau_c$ conservador.
